\section{Полугруппы операторов и инфинитезимальные генераторы}

Ранее мы установили, что стационарное распределение $\pi$ эргодической цепи Маркова удовлетворяет векторному уравнению $Q^T \pi = 0$, что эквивалентно инвариантности $\pi^T P(t) = \pi^T$ для любого $t \ge 0$. При этом матрица переходных вероятностей удовлетворяет полугрупповому свойству: $P(s+t) = P(s) P(t)$. Для дальнейшего анализа формализуем это наблюдение в терминах функционального анализа.

\begin{definition}{Полугруппа операторов}
    Пусть $X$ — банахово (линейное нормированное полное) пространство. Семейство ограниченных линейных операторов $\{P_t\}_{t \ge 0}$, действующих из $X$ в $X$, называется \textbf{полугруппой}, если выполнены следующие условия:
    \begin{enumerate}
        \item $P_{t+s} = P_t P_s$ для любых $s, t \ge 0$;
        \item $P_0 = I$, где $I$ — тождественный оператор.
    \end{enumerate}
    Обратимость операторов $P_t$ при этом не требуется.
\end{definition}

\begin{definition}{Непрерывность полугруппы}
    Полугруппа $\{P_t\}_{t \ge 0}$ называется \textbf{сильно непрерывной}, если для любого элемента $\varphi \in X$ выполнено условие:
    \[
        \lim_{h \to +0} \|P_h \varphi - \varphi\|_X = 0.
    \]
    Если сходимость выполняется по операторной норме в пространстве ограниченных операторов $\mathcal{B}(X)$, то есть $\lim\limits_{h \to +0} \|P_h - I\|_{\mathcal{B}(X)} = 0$, то полугруппа называется \textbf{равномерно непрерывной}.
\end{definition}

\begin{definition}{Инфинитезимальный генератор}
    Линейный оператор $G$, определяемый пределом
    \[
        G \varphi = \lim_{h \to +0} \frac{P_h \varphi - \varphi}{h},
    \]
    называется \textbf{генератором} полугруппы $\{P_t\}_{t \ge 0}$. Его областью определения является множество всех $\varphi \in X$, для которых указанный предел существует по норме пространства $X$.
\end{definition}

\begin{statement}
    Для равномерно непрерывной полугруппы генератор $G$ является ограниченным оператором, а сама полугруппа однозначно восстанавливается в виде экспоненциального ряда:
    \[
        P_t = e^{Gt}.
    \]
\end{statement}

\begin{note}
    Инфинитезимальный генератор тесно связан с абстрактной задачей Коши. Для дифференциального уравнения 
    \[
        \begin{cases} 
            \frac{d}{dt}u(t) = G u(t), \\ 
            u(0) = \varphi,
        \end{cases}
    \]
    где $u \colon [0, +\infty) \to X$, решение задается прямым действием полугруппы на начальное условие: $u(t) = P_t \varphi$.
\end{note}

Применим изложенный аппарат к марковским цепям. Пусть $S$ — не более чем счетное множество состояний марковской цепи $X(t)$, а $f \colon S \to \R$ — ограниченная функция состояний. Вычислим математическое ожидание $\E_{\vec{\alpha}} f(X(t))$ при заданном начальном распределении $\vec{\alpha} = (\alpha^{(1)}, \alpha^{(2)}, \dots)^T$:
\begin{align*}
    \E_{\vec{\alpha}} f(X(t)) &= \sum_{j \in S} \P(X(t)=j) f(j) \\
    &= \sum_{j \in S} \left( \sum_{i \in S} \P(X(t)=j \mid X(0)=i) \P(X(0)=i) \right) f(j) \\
    &= \sum_{i \in S} \alpha^{(i)} \left( \sum_{j \in S} p_{ij}(t) f(j) \right).
\end{align*}

Вводя вектор-столбец $\vec{f} = (f^{(1)}, f^{(2)}, \dots)^T \doteq (f(1), f(2), \dots)^T$ и вектор-столбец начального распределения $\vec{\alpha}$, перепишем результат через скалярное произведение:
\[
    \E_{\vec{\alpha}} f(X(t)) = \langle \vec{\alpha}, P(t) \vec{f} \rangle.
\]

Из прямого уравнения Колмогорова известно, что $\frac{d}{dt}P(t) = P(t) Q$. Интегрируя это матричное уравнение на отрезке $[0, t]$, получаем:
\[
    P(t) = I + \int\limits_0^t P(s) Q \, ds.
\]
Подставим полученное интегральное представление в выражение для математического ожидания:
\begin{align*}
    \langle \vec{\alpha}, P(t) \vec{f} \rangle &= \left\langle \vec{\alpha}, \left( I + \int\limits_0^t P(s) Q \, ds \right) \vec{f} \right\rangle \\
    &= \langle \vec{\alpha}, \vec{f} \rangle + \int\limits_0^t \langle \vec{\alpha}, P(s) (Q \vec{f}) \rangle \, ds.
\end{align*}
Заметим, что скалярное произведение $\langle \vec{\alpha}, \vec{f} \rangle$ совпадает с $\E_{\vec{\alpha}} f(X(0))$. Для интерпретации подынтегрального слагаемого рассмотрим вектор-столбец $Q \vec{f}$ как дискретное представление некоторой новой функции $Qf \colon S \to \R$, действующей на пространстве состояний:
\[
    (Qf)(i) = \sum_{j \in S} Q_{ij} f(j).
\]
При такой записи выражение $\langle \vec{\alpha}, P(s) (Q \vec{f}) \rangle$ приобретает смысл математического ожидания функции $Qf$ в момент времени $s$, то есть равно $\E_{\vec{\alpha}} [(Qf)(X(s))]$. Объединяя эти соотношения, приходим к следующему результату.

\begin{theorem}{Формула Дынкина}
    Для любой ограниченной функции $f$ и начального распределения $\vec{\alpha}$ справедливо соотношение:
    \[
        \E_{\vec{\alpha}} f(X(t)) = \E_{\vec{\alpha}} f(X(0)) + \int\limits_0^t \E_{\vec{\alpha}} [(Qf)(X(s))] \, ds,
    \]
    где $Qf$ — скалярная функция на пространстве состояний, полученная действием генератора $Q$ на исходную функцию $f$.
\end{theorem}

\begin{example}{Выражение для матрицы Q}
    Покажем, что инфинитезимальная матрица $Q$ действительно выступает в роли генератора полугруппы. 
\end{example}
\begin{eproof}
Пусть $\vec{\alpha} = \vec{e}_i$, где $\vec{e}_i$ — $i$-й базисный вектор. Это соответствует детерминированному начальному состоянию $X(0) = i$, откуда $\E_{\vec{e}_i} f(X(t)) = \vec{e}_i^T P(t) \vec{f}$. Вычислим действие оператора $Q$ на функцию $f$ в точке $i$:
    \begin{align*}
        (Qf)(i) &= \vec{e}_i^T (Q\vec{f}) = \vec{e}_i^T (P'(0) \vec{f}) \\
        &= \vec{e}_i^T \lim_{h \to +0} \frac{P(h)\vec{f} - P(0)\vec{f}}{h} \\
    &= \lim_{h \to +0} \frac{\vec{e}_i^T P(h)\vec{f} - \vec{e}_i^T P(0)\vec{f}}{h} \\
        &= \lim_{h \to +0} \frac{\E_{\vec{e}_i} f(X(h)) - f(i)}{h}.
    \end{align*}

    В качестве частного случая подставим в полученный предел индикаторную функцию фиксированного состояния $j \in S$: $f(s) = I_{\{s = j\}}$. Тогда $f(i) = \delta_{ij}$ (символ Кронекера), $\vec{f} = \vec{e}_j$, а $\E_{\vec{e}_i} f(X(h)) = \vec{e}_i^T P(t) \vec{f} = \vec{e}_i^T P(t) \vec{e}_j = p_{ij} = \P(X(h)=j \mid X(0)=i)$. 
    
    В зависимости от совпадения индексов $i$ и $j$ получаем два тождества:
    \begin{itemize}
        \item При $i \neq j$:
        \[
            Q_{ij} = \lim_{h \to +0} \frac{\P(X(h)=j \mid X(0)=i)}{h} = q_{ij}.
        \]
        \item При $i = j$:
        \[
            Q_{ii} = \lim_{h \to +0} \frac{\P(X(h)=i \mid X(0)=i) - 1}{h} = -q_i.
        \]
    \end{itemize}
    Полученные выражения в точности совпадают с классическими определениями интенсивностей переходов цепи Маркова непрерывного времени.
\end{eproof}

\begin{example}{Вычисление математического ожидания процесса}
    Рассмотрим марковский процесс $X(t)$, описывающий число заявок в системе в момент времени $t$. Пусть поступление новых заявок образует пуассоновский поток с интенсивностью $\lambda > 0$, а интенсивность обработки пропорциональна текущему числу заявок и составляет $\mu X(t)$. Требуется найти математическое ожидание $m(t) = \E X(t)$ при заданном начальном состоянии $X(0) = s_0$.
\end{example}
\begin{eproof}
    Найдем действие инфинитезимального генератора $Q$ на произвольную ограниченную функцию $f \colon S \to \R$. Вычислим математическое ожидание $\E_{\vec{e}_i} f(X(h))$, раскладывая вероятности переходов за малое время $h$ с точностью до $o(h)$:
    \begin{align*}
        (Qf)(i) &= \lim_{h \to +0} \frac{\E_{\vec{e}_i} f(X(h)) - f(i)}{h} \\
        &= \lim_{h \to +0} \frac{1}{h} \Big[ f(i+1) \P(X(h)=i+1 \mid X(0)=i) + f(i-1) \P(X(h)=i-1 \mid X(0)=i) \\
        &\quad + f(i) \P(X(h)=i \mid X(0)=i) - f(i) + o(h) \Big] \\
        &= \lim_{h \to +0} \frac{1}{h} \Big[ f(i+1) \lambda h + f(i-1) \mu i h + f(i) (1 - \mu i h - \lambda h) - f(i) + o(h) \Big].
    \end{align*}
    Сокращая на $h$ и переходя к пределу, получаем выражение для генератора:
    \[
        (Qf)(i) = \lambda (f(i+1) - f(i)) + \mu i (f(i-1) - f(i)).
    \]
    
    Применим формулу Дынкина к функции $f(s) = s$. Подставляя эту функцию в найденное выражение, находим значение $(Qf)$ в произвольном случайном состоянии $X(s)$:
    \[
        (Qf)(X(s)) = \lambda (X(s) + 1 - X(s)) + \mu X(s) (X(s) - 1 - X(s)) = \lambda - \mu X(s).
    \]
    
    Тогда по формуле Дынкина для математического ожидания получаем интегральное уравнение:
    \begin{align*}
        \E f(X(t)) &= \E X(0) + \int\limits_0^t \E [(Qf)(X(s))] \, ds \\
        m(t) &= s_0 + \int\limits_0^t \E [\lambda - \mu X(s)] \, ds \\
        &= s_0 + \int\limits_0^t (\lambda - \mu m(s)) \, ds.
    \end{align*}
    
    Дифференцируя обе части по $t$, переходим к классической задаче Коши для линейного дифференциального уравнения первого порядка:
    \[
        \begin{cases}
            m'(t) = \lambda - \mu m(t), \\
            m(0) = s_0.
        \end{cases}
    \]

    \newpage
    
    Решая это дифференциальное уравнение, окончательно находим математическое ожидание числа заявок в системе:
    \[
        m(t) = s_0 e^{-\mu t} + \frac{\lambda}{\mu} (1 - e^{-\mu t}).
    \]

    Проиллюстрируем характер сходимости переходного процесса к стационарному режиму. Положим интенсивность обслуживания $\mu = 1$, тогда отношение $\lambda / \mu$ совпадает с параметром нагрузки $\rho = \lambda$. Динамика математического ожидания $m(t)$ для различных начальных состояний $s \in \{0, 2, 5, 10\}$ и варьируемой нагрузки $\rho$ приведена на графиках ниже.
    
    \definecolor{ExampPurple}{RGB}{102, 51, 153}
    \begin{center}
    \pgfplotsset{
    local group style/.style={
        width=\linewidth, 
        height=5.5cm,
        xlabel={$t$},
        ylabel={$m(t)$},
        xmin=0, xmax=5,
        grid=major,
        domain=0:5,
        samples=100,
        thick,
        no markers,
        legend style={
            fill=Snow2,
            fill opacity=0.8,
            text opacity=1,
            draw=black!30
        },
        cycle list={
            {Red3, solid},
            {DeepSkyBlue4, dashed},
            {SpringGreen4, dashdotted},
            {ExampPurple, dotted},
            {Tan1, densely dashed}
        }
    }
}

\begin{minipage}{0.48\textwidth}
    \centering
    \begin{tikzpicture}
    \begin{axis}[
        local group style, 
        title={Начальное состояние $s_0 = 0$}, 
        legend pos=north east
    ]
        \addplot {1 - exp(-x)}; \addlegendentry{$\rho=1$}
        \addplot {3*(1 - exp(-x))}; \addlegendentry{$\rho=3$}
        \addplot {5*(1 - exp(-x))}; \addlegendentry{$\rho=5$}
        \addplot {7*(1 - exp(-x))}; \addlegendentry{$\rho=7$}
        \addplot {9*(1 - exp(-x))}; \addlegendentry{$\rho=9$}
    \end{axis}
    \end{tikzpicture}
\end{minipage}
\hfill
\begin{minipage}{0.48\textwidth}
    \centering
    \begin{tikzpicture}
    \begin{axis}[
        local group style, 
        title={Начальное состояние $s_0 = 2$}, 
        legend pos=north east
    ]
        \addplot {2*exp(-x) + 1 - exp(-x)}; \addlegendentry{$\rho=1$}
        \addplot {2*exp(-x) + 3*(1 - exp(-x))}; \addlegendentry{$\rho=3$}
        \addplot {2*exp(-x) + 5*(1 - exp(-x))}; \addlegendentry{$\rho=5$}
        \addplot {2*exp(-x) + 7*(1 - exp(-x))}; \addlegendentry{$\rho=7$}
        \addplot {2*exp(-x) + 9*(1 - exp(-x))}; \addlegendentry{$\rho=9$}
    \end{axis}
    \end{tikzpicture}
\end{minipage}

\vspace{1.5em}

\begin{minipage}{0.48\textwidth}
    \centering
    \begin{tikzpicture}
    \begin{axis}[
        local group style, 
        title={Начальное состояние $s_0 = 5$}, 
        legend pos=north east
    ]
        \addplot {5*exp(-x) + 1 - exp(-x)}; \addlegendentry{$\rho=1$}
        \addplot {5*exp(-x) + 3*(1 - exp(-x))}; \addlegendentry{$\rho=3$}
        \addplot {5*exp(-x) + 5*(1 - exp(-x))}; \addlegendentry{$\rho=5$}
        \addplot {5*exp(-x) + 7*(1 - exp(-x))}; \addlegendentry{$\rho=7$}
        \addplot {5*exp(-x) + 9*(1 - exp(-x))}; \addlegendentry{$\rho=9$}
    \end{axis}
    \end{tikzpicture}
\end{minipage}
\hfill
\begin{minipage}{0.48\textwidth}
    \centering
    \begin{tikzpicture}
    \begin{axis}[
        local group style, 
        title={Начальное состояние $s_0 = 10$}, 
        legend pos=north east
    ]
        \addplot {10*exp(-x) + 1 - exp(-x)}; \addlegendentry{$\rho=1$}
        \addplot {10*exp(-x) + 3*(1 - exp(-x))}; \addlegendentry{$\rho=3$}
        \addplot {10*exp(-x) + 5*(1 - exp(-x))}; \addlegendentry{$\rho=5$}
        \addplot {10*exp(-x) + 7*(1 - exp(-x))}; \addlegendentry{$\rho=7$}
        \addplot {10*exp(-x) + 9*(1 - exp(-x))}; \addlegendentry{$\rho=9$}
    \end{axis}
    \end{tikzpicture}
    \end{minipage}
    \end{center}
\end{eproof}

\begin{note}[1]
    При $t \to \infty$ математическое ожидание числа заявок стремится к величине $\lambda/\mu$, не зависящей от начального состояния $s_0$. Рассмотренный марковский процесс строго соответствует модели системы массового обслуживания $M/M/\infty$ с параметрами $\lambda$ и $\mu$, поскольку суммарная интенсивность обслуживания кратна числу заявок в системе. Найденный асимптотический предел в точности совпадает со средним числом заявок в стационарном режиме работы такой системы.

    Обратим внимание, что в классической теории массового обслуживания (по нотации Кендалла) система $M/M/\infty$ описывает математическую модель с бесконечным числом обслуживающих приборов. При этом не важно, конечна или бесконечна очередь. Поскольку число серверов ничем не ограничено, любая поступающая в систему заявка немедленно занимает свободный прибор и начинает обрабатываться.

    Классическим физическим примером такой модели выступает идеализированная система самообслуживания или автостоянка с бесконечным числом мест.
    
\end{note}

\begin{note}[2]
    Аппарат полугрупп операторов и вытекающая из него формула Дынкина позволяют анализировать переходные процессы в системе. Ранее использованные методы, основанные на решении алгебраического уравнения $Q^T \pi = 0$, давали исключительно стационарные (предельные) характеристики. Новый подход дает явные зависимости моментов распределения от времени.
\end{note}